\documentclass[11pt]{article}
\usepackage[a4paper,margin=2.5cm]{geometry}
\usepackage{amsmath}\usepackage{amssymb}\usepackage{hyperref}
\usepackage{microtype}
\setlength{\parskip}{6pt}
\setlength{\parindent}{0pt}

\begin{document}

\begin{flushright}
2026-05-31
\end{flushright}

\noindent Prof.\ László T.\ Kóczy

\textit{Subject}: A learnable Kóczy fuzzy signature as a
differentiable hypergraph layer — early results and a request for
your feedback.

Dear Professor Kóczy,

I have been working on an explicit, learnable realisation of fuzzy
signatures as a differentiable neural-network layer, and I would
like to share the early results with you. The motivation is your
classical construction: a fuzzy signature is a hierarchical tree of
fuzzy concepts in which internal nodes are aggregation operators
(t-norms, t-conorms, OWA) and leaves carry elementary memberships.
Almost everything in the construction is interpretable; only the
membership functions and the aggregator weights are tuned in the
classical setting.

The construction I have built keeps the structure but makes
\emph{every} learnable element a named fuzzy-systems primitive:

\begin{itemize}
  \item Each per-vertex feature is a learnable Catmull--Rom
        membership function $\mu_\theta\colon X\to[0,1]$.
  \item The signed branches of the architecture are an Atanassov
        intuitionistic-fuzzy pair $(\mu^+, \mu^-)$ with hesitancy
        $\pi = 1 - \mu^+ - \mu^-$.
  \item The polarity gate that mixes the pair is a Zadeh sigmoidal
        hedge with learnable dilation $\tau$ per channel.
  \item Vertex-to-edge aggregation is a categorical t-norm ($T_G$ /
        $T_P$ / $T_L$); edge-to-vertex is a categorical t-conorm.
  \item Multi-arity (multi-scale) fusion across receptive-field
        sizes is an OWA with $\alpha_k$ softmax weights.
  \item The output head is a Sugeno weighted-average defuzzification
        followed by the TSK output rule.
\end{itemize}

I attach a 20-page mathematical background that lays out this
correspondence in full, along with the architectural plan and a
robotic-collaboration application story.

Empirically, the architecture has been validated on MNIST at a small
parameter budget (test\_acc $0.5354$ at 9{,}642 parameters; the
$\alpha_k$ multi-scale mixer emergently specialises so that the
first layer prefers the fine receptive field and the last layer the
coarse one — exactly the multi-scale story HSiKAN-vision relies on,
re-emerging in the fuzzy-signature setting). The architecture also
extends to a pose-detection variant; the validated single-frame
synthetic-pose model is now the basis for a planned two-robot
collaborative-behavior test bed.

I emphasise the honest current limits: the MNIST result is
single-seed; the architecture is not yet at HSiKAN parity ($\sim 0.95$+);
the pose-detection comparison against simpler baselines is in flight.
The point of this letter is to ask for your feedback on the
direction, not to claim a final result. In particular:

\begin{itemize}
  \item Does the named-primitive correspondence preserve what you
        consider essential about fuzzy signatures?
  \item Is the Atanassov-pair reading of the signed-branch
        construction the right inductive bias for an
        \emph{intuitionistic} fuzzy signature, or do you see a
        different formulation that better captures hesitancy?
  \item For the robotic-collaboration story (two-robot test bed,
        cooperative-handover / competitive-grasp / independent-parallel
        interactions), does the framework's hierarchical aggregation
        align with how you envision fuzzy signatures supporting
        human-robot collaboration?
\end{itemize}

I would value any feedback you can offer, and I am happy to share
the full implementation, the test matrix, or any other technical
detail you would like to see.

With deep respect for your foundational work,

\bigskip

\noindent\textsc{[Sender]}

\bigskip

\noindent\textbf{Attachments:}
\begin{enumerate}
  \item \texttt{background.pdf} — Mathematical Background for the
        Learnable Kóczy Fuzzy Signature Layer (20 pp). Atanassov
        IFS, classical t-norms / t-conorms / negations, Kóczy
        fuzzy signatures, defuzzification methods, Zadeh
        linguistic hedges, OWA aggregators, the full HSiKAN-Kóczy
        correspondence table, $[0,1]$-preservation theorem,
        TSK-output-rule derivation, worked examples.
  \item \texttt{plan\_fuzzy\_signature.pdf} — Architectural plan
        for the FuzzySignatureLayer (7 pp).
  \item \texttt{plan\_robotics.pdf} — Robotics / behavior /
        collaboration application plan with the two-robot test
        case (7 pp).
\end{enumerate}

\end{document}
